% Insert after the fixed-bank replay subsection and before the exemplar bridge.
\subsection{When replay reverses concentration, and how fast}
\label{app:theory-replay-finite-horizon}

The restoring logit field $\rho(\bar w-p)$ has a shared foundation with
inverse probability scaling: the ideal IPS field of
\citet[Theorem~4.1]{sinha2026expected}, normalized to a fixed target $w$, is
$\bar w-p$. The uniform-correct target and its forward conditional-KL
formulation are also established by \citet{lochab2026ucpo}. Here the target is
restricted to a discovered bank and evaluated from stored exemplars. The
question below is how its restoring field competes with the fresh binary
objective over a finite interval. Neither a uniform target nor its
cross-entropy gradient is a new contribution of this analysis.

Fix a bank $\mathcal B\subseteq\mathcal C$ of $k\ge2$ modes, with uniform
categorical weights, and consider Equation~\eqref{eq:replay-mean-flow}.
Write $c(P)$ for its fresh coefficient and define
\begin{equation}
 \begin{gathered}
 M_{\mathcal B}=\sum_{b\in\mathcal B}p_b,\qquad
 r_b=\frac{p_b}{M_{\mathcal B}},\qquad
 C_{\mathcal B}=\sum_{b\in\mathcal B}r_b^2,\\
 a_\rho(t)=\rho-c(P(t))(1-P(t)).
 \end{gathered}
 \label{eq:replay-finite-notation}
\end{equation}
Thus $1-C_{\mathcal B}$ is \pmd{} conditional on the bank. It equals the
paper's success-conditional \pmd{} when $\mathcal B=\mathcal C$; otherwise
it measures only allocation inside the bank.

\begin{theorem}[Exact replay threshold for bank diversity]
\label{thm:replay-finite-threshold}
Under the categorical mean-flow assumptions and a fixed uniform bank,
\begin{align}
 \frac{d}{dt}\log\frac{p_i}{p_j}
  &=-a_\rho M_{\mathcal B}(r_i-r_j),
   &&i,j\in\mathcal B,\label{eq:replay-finite-ratio}\\
 \dot r_b&=-a_\rho M_{\mathcal B}r_b(r_b-C_{\mathcal B}),
   &&b\in\mathcal B,\label{eq:replay-finite-conditional}\\
 \frac{d}{dt}(1-C_{\mathcal B})
  &=2a_\rho M_{\mathcal B}
    \left(\sum_{b\in\mathcal B}r_b^3-C_{\mathcal B}^2\right).
 \label{eq:replay-finite-pcmd}
\end{align}
For a nonuniform $r$, bank-conditional diversity strictly increases when
$\rho>c(P)(1-P)$, strictly decreases when $\rho<c(P)(1-P)$, and has zero
instantaneous derivative at equality. These conclusions require no coverage
of correct modes outside the bank.
\end{theorem}

\begin{proof}
For $b\in\mathcal B$, the score identities give
$\dot z_b=\rho/k-a_\rho p_b$. Subtracting two such equations gives
Equation~\eqref{eq:replay-finite-ratio}. The distribution $r$ is the softmax
of the logits restricted to $\mathcal B$, so
$\dot r_b=r_b(\dot z_b-\sum_jr_j\dot z_j)$, proving
Equation~\eqref{eq:replay-finite-conditional}. Differentiating
$C_{\mathcal B}$ gives Equation~\eqref{eq:replay-finite-pcmd}.
Its parenthesis is $\operatorname{Var}_{b\sim r}(r_b)$, which is strictly
positive for a nonuniform distribution with positive coordinates.
\end{proof}

A positive dose therefore need not improve diversity immediately.
The thresholds for Dr.GRPO and the specified MaxRL estimator are, respectively,
\[
 \rho>\frac{G-1}{G}(1-P),\qquad
 \rho>\frac{(1-P)[1-(1-P)^{G-1}]}{P}.
\]
The asymptotic result for every $\rho>0$ and this finite-time condition are
compatible: the fresh concentration pressure vanishes as $P\to1$.

\begin{theorem}[Finite recovery from accumulated replay exposure]
\label{thm:replay-finite-recovery}
In the setting of Theorem~\ref{thm:replay-finite-threshold}, suppose
$a_\rho(u)\ge0$ for $u\in[T,t]$, and define
\[
 S(T,t)=\int_T^t a_\rho(u)M_{\mathcal B}(u)\,du,
 \qquad D(t)=\log\frac{\max_{b\in\mathcal B}r_b(t)}
                            {\min_{b\in\mathcal B}r_b(t)}.
\]
Then
\begin{equation}
 e^{D(t)}-1\le(e^{D(T)}-1)e^{-S(T,t)/k},
 \qquad
 \min_b r_b(t)\ge\frac{1}{1+(k-1)e^{D(t)}}.
 \label{eq:replay-finite-recovery}
\end{equation}
For $0<\delta<1/k$, put
$D_\delta=\log[(1/\delta-1)/(k-1)]$. If $D(T)>D_\delta$, the condition
\begin{equation}
 S(T,t)\ge k\log\frac{e^{D(T)}-1}{e^{D_\delta}-1}
 \label{eq:replay-finite-visibility-exposure}
\end{equation}
is sufficient for every $r_b(t)\ge\delta$. If $D(T)\le D_\delta$, this
conclusion already holds and is preserved on the interval. In particular,
when $a_\rho\ge a_0>0$ and $M_{\mathcal B}\ge M_0>0$ on the interval,
its required duration is at most the right-hand side of
Equation~\eqref{eq:replay-finite-visibility-exposure} divided by $a_0M_0$.
At that endpoint every unconditional bank probability is at least
$M_0\delta$.
\end{theorem}

\begin{proof}
Bank-coordinate ordering is preserved by
Equation~\eqref{eq:replay-finite-ratio}: a difference of two bank logits
satisfies a scalar equation proportional to that difference, and ties persist.
Writing $r_{\max}$ and $r_{\min}$ for the extremal coordinates,
\[
 \dot D=-a_\rho M_{\mathcal B}(r_{\max}-r_{\min}),
 \qquad
 r_{\max}-r_{\min}
   =r_{\max}(1-e^{-D})\ge\frac{1-e^{-D}}{k}.
\]
Thus $\frac{d}{dt}(e^D-1)\le
 -(a_\rho M_{\mathcal B}/k)(e^D-1)$; integration proves the first bound.
Since $1\le r_{\min}+(k-1)r_{\max}$, the second follows. Substituting
$D_\delta$ and solving the first bound for the exposure proves the remaining
claims. When $D(T)=0$, uniformity is preserved directly.
\end{proof}

The bound specifies the interval on which restoration occurs. It does not
assume that bank mass is monotone or that correctness stays above its value
at $T$. Both enter through the displayed exposure; a bound on unconditional
visibility additionally needs the stated bank-mass lower bound.

\begin{proposition}[Discrete exact-gradient contraction]
\label{prop:replay-discrete-contraction}
Replace the continuous update by simultaneous exact mean-gradient steps,
while retaining independent logits and a fixed uniform bank:
\[
 z_{n+1}=z_n+\eta_n\{c(P_n)\nabla_zP_n+
                         \rho(\bar w-p_n)\},\qquad \eta_n\ge0.
\]
Define $D_n$ as above and
$\gamma_n=\eta_n M_{\mathcal B,n}
 [\rho-c(P_n)(1-P_n)]$.
If $0\le\gamma_n\le2$, the order of all bank probabilities is preserved
at this step and
\begin{equation}
 e^{D_{n+1}}-1\le
 \left(1-\frac{1-e^{-\gamma_n}}{k}\right)(e^{D_n}-1).
 \label{eq:replay-discrete-contraction}
\end{equation}
Consequently, if $0\le\gamma_n\le1$ for $n=0,\ldots,N-1$, then
\[
 e^{D_N}-1\le(e^{D_0}-1)
  \exp\!\left(-\frac{1-e^{-1}}{k}
                    \sum_{n=0}^{N-1}\gamma_n\right).
\]
\end{proposition}

\begin{proof}
For two bank logits with difference $d=z_i-z_j\ge0$,
\[
 d_{\rm next}=d-\gamma_n(r_i-r_j),\qquad
 r_i-r_j=(r_i+r_j)\tanh(d/2)\le d/2.
\]
Hence $d_{\rm next}\ge0$ when $\gamma_n\le2$. The same extremal indices
therefore give
$D_{n+1}=D_n-\gamma_n(r_{\max}-r_{\min})$.
For $x\in[0,1]$, convexity of $e^{-\gamma_nx}$ yields
$1-e^{-\gamma_nx}\ge(1-e^{-\gamma_n})x$. Set
$x=r_{\max}-r_{\min}$ and use
$e^{D_n}x\ge(e^{D_n}-1)/k$ from the preceding proof. This gives
Equation~\eqref{eq:replay-discrete-contraction}. Multiplication over steps,
$1-y\le e^{-y}$, and
$1-e^{-\gamma}\ge(1-e^{-1})\gamma$ for $0\le\gamma\le1$ prove the last
bound.
\end{proof}

These discrete statements concern exact mean gradients, without clipping,
sampling noise, optimizer state or alternating replay visits. They establish
a finite-step result for the specified update, rather than a guarantee for
AdamW. Uniform categorical weights are also essential to the stated target.
For fixed nonuniform weights $w_b$, the exact conditional equation is instead
\begin{equation}
 \dot r_b=r_b\left\{
   \rho\left(w_b-\sum_jr_jw_j\right)
   -a_\rho M_{\mathcal B}(r_b-C_{\mathcal B})\right\}.
 \label{eq:replay-finite-weighted}
\end{equation}
The inverse-length weights induced by the implemented mean-token loss belong
to this weighted case. The uniform contraction theorem therefore describes
an equal-length categorical specialization, and does not assert uniform key
probabilities for neural exemplar replay.
